\documentclass[10pt]{beamer}
\usetheme{default}
\setbeamercovered{invisible}
\setbeamertemplate{navigation symbols}{}
\setbeamertemplate{footline}{
    \flushright{\hfill \insertframenumber{}/\inserttotalframenumber}
}

\usepackage{listings}

% User-defined colors.
\definecolor{DarkGreen}{rgb}{0, .5, 0}
\definecolor{DarkBlue}{rgb}{0, 0, .5}
\definecolor{DarkRed}{rgb}{.5, 0, 0}
\definecolor{LightGray}{rgb}{.95, .95, .95}
\definecolor{White}{rgb}{1.0,1.0,1.0}
\definecolor{darkblue}{rgb}{0,0,0.9}
\definecolor{darkred}{rgb}{0.8,0,0}
\definecolor{darkgreen}{rgb}{0.0,0.85,0}

% Settings for listing class.
\lstset{
  language=C++,                        % The default language
  basicstyle=\small\ttfamily,          % The basic style
  backgroundcolor=\color{White},       % Set listing background
  keywordstyle=\color{DarkBlue}\bfseries, % Set keyword style
  commentstyle=\color{DarkGreen}\itshape, % Set comment style
  stringstyle=\color{DarkRed}, % Set string constant style
  extendedchars=true % Allow extended characters
  breaklines=true,
  basewidth={0.5em,0.4em},
  fontadjust=true,
  linewidth=\textwidth,
  breakatwhitespace=true,
  showstringspaces=false,
  lineskip=0ex, %  frame=single
}

\begin{document}
    \title{A complete Newton solver\protect\\ using \texttt{Eigen}}
    \author{Paolo Joseph Baioni}
    \date{\today}

\begin{frame}[plain, noframenumbering]
    \maketitle
\end{frame}

\begin{frame}{Newton solver}
	This example (inspired by \texttt{Examples/src/NewtonSolver}) presents a set of tools implementing Newton methods for non-linear systems, using the \texttt{Eigen} library.
	
	The code structure is the following:
	\begin{itemize}\small
		\item \texttt{NewtonTraits} defines the common types used by the solver classes, ensuring a uniform interface.
		\item \texttt{JacobianBase} is an abstract base class representing the action of a Jacobian (or approximate Jacobian).
		\item \texttt{FullJacobian} uses a user-provided Jacobian matrix.
		\item \texttt{DiscreteJacobian} approximates the Jacobian by finite differences.
		\item \texttt{JacobianFactory} returns a concrete object in the \texttt{JacobianBase} family.
		\item \texttt{Newton} applies the Newton iteration given a non-linear system and a Jacobian object.
		\item \texttt{NewtonOptions} and \texttt{NewtonResult} collect input parameters and output data.
	\end{itemize}
\end{frame}

\begin{frame}{Exercise}
	Consider the problem
	\[
	\mathbf{f}(x, y) =
	\begin{bmatrix}
		(x-1)^2 + 0.1(y - 5)^2 \\
		1.5 -x - 0.1y
	\end{bmatrix}
	=
	\mathbf{0}.
	\]
	
	Starting from the provided solution sketch:
	\begin{enumerate}\small
		\item Implement the \texttt{NewtonTraits} class, defining the common types used by the solver.
		\item Implement the \texttt{FullJacobian} class (derived from \texttt{JacobianBase}), which, given the full Jacobian matrix, solves the linear system by direct \(LU\) factorization.
		\item Implement the \texttt{DiscreteJacobian} class (derived from \texttt{JacobianBase}), which approximates the Jacobian by finite differences and solves the resulting linear system by direct \(LU\) factorization.
		\item Implement a \texttt{JacobianFactory} function returning an instance of \texttt{FullJacobian} or \texttt{DiscreteJacobian}, depending on the selected Jacobian type.
		\item Solve the problem above using both the full and the discrete approach.
		\item Extra: generalize the problem type (scalar, vector, ...) with templates, while keeping runtime Jacobian selection.
	\end{enumerate}
\end{frame}

\end{document}
